% Minimal retained preamble for lane-p2 cards taken from
% sources/hitting-sets-hypergraphs-2307.12154v3/source0.tex (arXiv:2307.12154v3).
%
% Copied verbatim from the author's own file:
%   line 3      : \documentclass[11pt,oneside]{amsart}
%   lines 43-49 : the author's macro definitions
%   lines 53-59 : the author's theorem declarations
% The base packages amsmath, amsfonts and amssymb are supplied here because amsart
% does not load amssymb by itself (LANE_BRIEF 18.1).
%
% Omitted author preamble lines and why (nothing was renamed or redefined):
%   line 10 (\usepackage[utf8]{inputenc}), line 11 (fontenc), line 12 (lmodern):
%       input encodings and font encodings fight the forced xelatex engine;
%   line 13 (geometry), line 34 (\baselinestretch), line 29 (parskip): pure layout,
%       and the assembler supplies the page setup;
%   lines 14-24 and 28 and 32 (array, fancyhdr, graphicx, float, caption, pgf, tikz,
%       tikzlings-addons, tikzlings-penguins, subcaption, adjustbox, verbatim): none of
%       the retained spans draws or floats anything, and the tikzlings package is not
%       part of the host TeX Live;
%   line 62 (\freefootnote): the retained spans contain no footnote.
\documentclass[11pt,oneside]{amsart}
\usepackage{amsmath}
\usepackage{amsfonts}
\usepackage{amssymb}

\newcommand{\R}{\mathbb{R}}
\newcommand{\N}{\mathbb{N}}
\newcommand{\HH}{\mathcal{H}}
\newcommand{\RR}{\mathcal{R}}
\newcommand{\A}{\mathcal{A}}
\newcommand{\B}{\mathcal{B}}
\newcommand{\est}{\emptyset}

\newtheorem{theorem}{Theorem}
\newtheorem{proposition}{Proposition}
\newtheorem{corollary}{Corollary}
\newtheorem{lemma}{Lemma}
\newtheorem{definition}{Definition}
\newtheorem{jel}{Notation}
\newtheorem{conj}{Conjecture}
\title{Shallow hitting sets in a hereditary hypergraph family bound the polychromatic threshold}
\author{Balázs Bursics, Bence Csonka, Luca Szepessy}
\date{arxiv:2307.12154v3; CC BY 4.0}
\begin{document}
\maketitle
\noindent\textit{Editorial note.} The counted unit is the article's own Lemma prop1 at source lines 151-155 together with the proof environment at source lines 369-374, whose header names it: Proof of Lemma prop1. The pairing is therefore non-adjacent (the article collects the proofs of its introduction in a later section), and it rests on the author's own named header rather than on adjacency. The statement contains the author's commented-out variant of its second sentence at source line 153, retained exactly as he left it.\par\medskip
\noindent\textit{Editorial note.} The retained context at source lines 120, 123-124, 128-129, 138-140 and 142-144 is the article's own definition of a polychromatic k-colouring, its notation for truncated hypergraphs and families, its definition of the parameter m\_H(k), its definition of a c-shallow hitting set and its definition of an induced subhypergraph. The surrounding paragraphs that carry citations are deliberately not retained, so the material carries no bibliography.\par\medskip
\noindent\textit{Editorial note.} The article is typeset with amsart plus an input encoding, a font encoding, geometry, array, fancyhdr, graphicx, float, caption, pgf, tikz, two tikzlings packages, subcaption, parskip, adjustbox and verbatim. The retained preamble keeps the author's class line and his own macro and theorem declarations and drops the rest: the encodings and fonts fight the forced xelatex engine, the layout packages are supplied by the assembler, and the retained spans draw nothing and contain no floats, so the drawing packages are not needed and the tikzlings family is not part of the host TeX Live. amsmath, amsfonts and amssymb are loaded explicitly because amsart does not load amssymb by itself. Nothing was renamed, redefined or added.\par\medskip
\noindent\textit{Editorial note.} The proof as retained is complete: the five sentences shrink the edges, consume one shallow hitting set per colour class, and observe that one uncoloured vertex remains in every edge after k minus one rounds.\par\medskip
\noindent\textit{Allowed standard prerequisites (editorial).}\begin{itemize}
\item Polychromatic colourings and the parameter m\_H(k). A polychromatic k-colouring of a hypergraph colours its vertices with k colours so that every edge meets every colour class; for a hypergraph family H the parameter m\_H(k) is the smallest threshold m such that every member admits a polychromatic k-colouring of its subhypergraph of edges of size at least m, and it is infinite when no such threshold exists; the article's own definitions are retained at source lines 120 and 128-129.. The counted lemma bounds m\_H(k) by c times k minus one plus one, so both the colouring notion and the parameter must be available.
\item Truncated hypergraphs and the hereditary assumption of the family. H at least m is obtained by deleting every edge of size below m and H exactly m keeps only the edges of size exactly m, with the analogous notation for a family; the counted lemma assumes that the family is closed under induced subhypergraphs and under shrinking one edge by one of its vertices; retained at source lines 123-124 and 142-144.. The proof repeatedly removes the vertices of a shallow hitting set and then shrinks edges, which is exactly what these two closure assumptions license.
\item Shallow hitting sets. For a positive integer c, a vertex set U is a c-shallow hitting set of a hypergraph when every edge meets U in at least one and at most c vertices; retained at source lines 138-140.. The counted lemma assumes the existence of a c-shallow hitting set for every member of the truncated family, and the proof consumes one such set per colour class.
\end{itemize}\medskip
A {\it polychromatic k-coloring} of a hypergraph $H$ is a $k$-coloring of its vertex set such that every hyperedge contains a vertex of all $k$ color classes. Observe that a polychromatic 2-coloring is the same as the usual proper 2-coloring of hypergraphs, where we require that no edge is monochromatic. By merging $j$ color classes of a polychromatic $k$-coloring we get a polychromatic $(k-j+1)$-coloring, so the condition of monochromatic $k$-colorability becomes stricter as $k$ increases. A trivial necessary condition for the existence of a polychromatic $k$-coloring is that all edges of $H$ must be of size at least $k$.

\begin{definition} Denote by $H_{\ge m}$ or simply $H_m$ the hypergraph obtained from $H$ by deleting all hyperedges of size smaller than $m$, and denote by $H_{=m}$ the hypergraph consisting of the hyperedges of $H$ with size exactly $m$. Similarly, for a hypergraph family $\mathcal{H}$ let $\HH_m=\HH_{\ge m}=\{H_{\ge m}:H\in\HH\}$ and $\HH_{=m}=\{H_{=m}:H\in\HH\}$.
\end{definition}

\begin{definition} Let $\mathcal{H}$ be a hypergraph family. Denote by $m_\mathcal{H}(k)$ the smallest positive integer $m$ such that for every $H\in\mathcal{H}$ there exists a polychromatic $k$-coloring of the hypergraph $H_{\ge m}$. If there is no such $m$, set $m_{\mathcal{H}}(k)=\infty$.
\end{definition}

\begin{definition}
 Let $H=(V,E)$ be a hypergraph, a $U\subset V$ vertex set $U\subset V$ is a {\it c-shallow hitting set} for a positive integer $c$ if for every hyperedge $e\in E$: $1\le|e\cap U|\le c.$
\end{definition} 

\begin{definition}
 Let $H=(V,E)$ be a hypergraph, and $V'\subset V$. Define $E'=\{e'\subset V'| \exists e\in E: e'=e\cap E'\}$, we say that $E'=(V',E')$ is the induced subhypergraph of $H$ on $V'$.
\end{definition}

\begin{lemma}\label{prop1}
Suppose that the hypergraph family $\mathcal{H}$ satisfies the condition that for arbitrary $H=(V,E)\in\HH$ and $V'\subset V$, the induced subhypergraph of $H$ on $V'$ is an element of $\HH$, and that every hyperedge $e\in E$ contains a vertex $v\in e$ such that for $e'=e\setminus \{v\}$ and $E'=E\setminus \{e\}\cup\{e'\}$ we have $H'=(V,E')\in\HH.$
%replacing $e$ with $e'=e\setminus v$ results in an element of $\HH$. 
Suppose that for every $m\ge c$ and every element of $\mathcal{H}_{=m}$ there exists a $c$-shallow hitting set. Then $m_\mathcal{H}(k) \leq c \cdot (k-1) +1$.
\end{lemma}

\begin{proof} [Proof of Lemma \ref{prop1}.] 
Let $H\in\HH$ be arbitrary, we need to show that $H_{\ge c(k-1)+1}$ can be colored polychromatically with $k$ colors.

First shrink all edges of $H_{\ge c(k-1)+1}$ to be of size exactly $c(k-1)+1$ in such a way that the resulting graph is in $\HH$, we can do this by our assumptions.
Then choose a $c$-shallow hitting set from our hypergraph, color its points to one fixed color, and leave out these vertices from the graph. The resulting graph is still an element of the hypergraph family, and all of its edges can be shrinked to contain exactly $c(k-2)+1$ vertices, and so that we get a graph from $\HH_{=c(k-2)+1}$. Then take another $c$-shallow hitting set, color it to the second color, and so on. After determining this way the first $k-1$ color classes, we still have at least one remaining vertex in every edge, color the remaining vertices to the $k$th color to get a desired $k$-coloring.
\end{proof}

\end{document}
